% Tabla compuesta: las cinco afirmaciones que vertebran el documento, cada una
% con la cifra que la sostiene, el artefacto donde verificarla y su matiz.
\begin{tabularx}{\textwidth}{>{\raggedright\arraybackslash}p{0.24\textwidth}lX}
\toprule
Afirmación & Cifra & Matiz que la acompaña \\
\midrule
El sistema aprende a ponderar a sus especialistas &
$+52\,\%$ &
El meta pasa de reparto uniforme a identificar a \textit{risk} en dos años, pero queda pegado a su
tope de 0,50 desde 2019 \\
\addlinespace
Lo aprendido es señal y no ruido &
$\rankic = 0{,}1004$ &
Frente a 0,0130 del mejor \textit{baseline}, pero 117 cohortes equivalen a $\approx 9$
observaciones independientes \\
\addlinespace
No es suerte &
7 de 8 contrastes &
Permutación $p=0{,}0001$ y placebos en cero; el Deflated Sharpe 0,930 no alcanza su umbral y se
reporta igualmente \\
\addlinespace
Bate al índice en la era reservada &
$+14{,}21\,\%$ &
Exceso geométrico con IR 0,959 y 2 de 2 años, pero el Rank-IC de esa era es indeterminado por falta
de cohortes cerradas \\
\addlinespace
La mejor forma de usar la señal es no interferirla &
$\text{TC}=0{,}247$ &
El perfil que no reordena gana; seis de siete perfiles sesgados destruyen alfa. La cartera
materializa una cuarta parte de la señal \\
\bottomrule
\end{tabularx}
